% =============================================================================
\section{Conclusion et perspectives}
% =============================================================================

\begin{frame}{Contributions et perspectives}
  \footnotesize
  \begin{columns}[T,onlytextwidth]
    \begin{column}{0.5\textwidth}
      \begin{exampleblock}{Contributions}
        \begin{itemize}
          \item Approche \textbf{conjointe} bloc + lits + ambulatoire, en
            dynamique.
          \item \textbf{Dix méthodes} comparées, validées à l'optimum exact.
          \item \textbf{Aléas} : plannings alternatifs (Date A / Date B) et
            adaptation stable.
          \item Écosystème complet : solveur, simulation Mesa, GUI, CLI.
        \end{itemize}
      \end{exampleblock}
    \end{column}
    \begin{column}{0.48\textwidth}
      \begin{alertblock}{Suite du travail}
        \begin{itemize}
          \item Valider le \textbf{proxy de spécialité} et les poids du coût.
          \item Intégrer l'\textbf{ambulatoire} dans le simulateur Mesa.
          \item Boucle journalière d'indicateurs opérationnels.
          \item Budgets plus longs pour départager le haut du tableau.
        \end{itemize}
      \end{alertblock}
    \end{column}
  \end{columns}

  \vspace{0.15cm}
  \begin{block}{Message clé}
    \centering
    L'originalité du fil rouge est la \textbf{gestion conjointe et coordonnée
    du bloc, des lits et de l'ambulatoire} dans un contexte dynamique.
  \end{block}
\end{frame}

% =============================================================================
% BIBLIOGRAPHIE
% =============================================================================
